% preamble
\documentclass[11pt,a4paper,twocolumn]{article}
\usepackage[norsk]{babel}
\usepackage[a4paper,
            top=2cm,
            bottom=2cm,
            left=3cm,
            right=3cm,
            marginparwidth=1.75cm]{geometry}
\usepackage{amsmath}
\usepackage[norsk]{cleveref}
\usepackage[shortlabels]{enumitem}
\usepackage{graphicx}
\usepackage[T1]{fontenc}
\usepackage[hyphens]{url}
\usepackage{lipsum} 
\usepackage{siunitx}
\sisetup{
  output-complex-root = \ensuremath{\mathit{j}},
  complex-root-position = before-number,
  exponent-product = \cdot,
  output-decimal-marker = {,}
}
\usepackage{tikz}
\usepackage{circuitikz}
\usetikzlibrary{arrows.meta,positioning,calc}
\usepackage{pgfplots}
\pgfplotsset{compat=1.18}
\setlength{\parindent}{0pt}

% metadata
\title{Rapportskriving i ELI1300 - Redusert mal}
\author{Odin Heggvold Bekkelund\\For Gruppe ADK57}
\date{31.08.2026} 

% definer en url til senere?
\urldef{\myoverleafurl}\url{https://tinyurl.com/2yfbtdez}

% selve dokumentet starter
\begin{document}
\maketitle % Setter sammen \title og \author og \date

\begin{abstract}
Sammendrag
\end{abstract}

% Innholdsfortegnelse automatisk generert:
%\tableofcontents

\section{Innledning}
Denne øvelsen vil omhandle å konvertere digitale signaler til en analog spenning såkalt digital til analog omformer (DA-omformer). Kretsen som skal praktisk bygges opp skal inneholde en DA omformer som bruker 3 bit. 

\subsection{Utstyrsliste}

1 2 3

\section{Teoretisk del}
I denne delen presenteres teori og sammenhenger som er viktige for å teste ulike DA kretser.

\subsection{Binærtall og antall digitale nivåer}

For et digitalt system med \(N\) bits er antall mulige binære kombinasjoner gitt ved:

\[
2^N.
\]

3 bit gir derfor \(2^3 = 8\) binære kombinasjoner.

\subsection{Fra digital verdi til analog utgang}

Sammenhengen mellom de sentrale delene i kretsen kan illustreres som:
\[
\begin{aligned}
	&\text{digitale spenningsnivåer} \\
	&\downarrow \\
	&\text{R--2R-nettverk} \\
	&\downarrow \\
	&\text{nodeanalyse} \\
	&\downarrow \\
	&v_{\mathrm{ut}}(v_2,v_1,v_0) \\
	&\downarrow \\
	&\text{DAC-karakteristikk} \\
	&\downarrow \\
	&\text{praktisk krets} \\
	&\downarrow \\
	&\text{belastning med LED}
\end{aligned}
\]

Utgangsspenningen kan dermed beskrives som en funksjon av de tre digitale inngangene:

\[
v_{\mathrm{ut}} = f(v_2,v_1,v_0).
\]

\subsection{Ohms lov og Kirchhoffs strømlov}

Ohms lov beskriver sammenhengen mellom spenning, strøm og motstand:

\[
i = \frac{v}{R}.
\]

Ved analyse av nodene i R--2R-nettverket benyttes Kirchhoffs strømlov (KCL). 
Summen av strømmene inn og ut av en node er lik null. For en node med spenning
\(v_x\) kan dette eksempelvis skrives som

\[
\frac{v_x-v_a}{R_1}
+
\frac{v_x-v_b}{R_2}
+
\frac{v_x-v_c}{R_3}
= 0.
\]


\subsection{Nodespenningsmetoden i R-2R-nettverket}
\begin{equation}
\begin{gathered}
I: \frac{v_a-v_0}{2R}+\frac{v_a-v_b}{R}+\frac{v_a}{2R} = 0 \\
II: \frac{v_b-v_a}{R} + \frac{v_b-v_1}{2R}+\frac{v_b-v_c}{R} = 0 \\
III: \frac{v_c-v_b}{R} + \frac{v_c-v_2}{2R} + \frac{v_c-v_{ut}}{R} = 0 \\
IV: \frac{v_{ut}-v_c}{R} + \frac{v_{ut}}{R_L} = 0 \\
\end{gathered}
\label{eqs:kcl}
\end{equation}

\Cref{eqs:kcl} viser oppsett av nodespenningsmetoden i kretsen i \cref{fig:r2r-krets}. Ligningssettet gir oss to utrykk for $v_{ut}$:

\[
\begin{gathered}
III': v_{ut} = 2v_c - \frac{v_0}{8} - \frac{v_1}{4} - \frac{v_2}{2} \\
\text{og}\\
IV': v_{ut} = \frac{R_L}{R + R_L}v_c.\\
\end{gathered}
\]

Satt sammen blir dette:
\[
\frac{R_L}{R + R_L}v_c = 2v_c - \frac{v_0}{8} - \frac{v_1}{4} - \frac{v_2}{2}\\.
\]

Som blir:
\begin{equation}
v_{ut} = \frac{R_L}{2R + R_L}(\frac{v_0}{8} + \frac{v_1}{4} + \frac{v_2}{2})\\
\label{eq:vutfromvin_full}
\end{equation}

\Cref{eq:vutfromvin_full} kan videre forenkles til \cref{eq:vutfromvin_short} når $R_L \to \infty$. Da går $\frac{R_L}{2R + R_L} \to 1$.

\begin{equation}
v_{ut} = \frac{v_0}{8} + \frac{v_1}{4} + \frac{v_2}{2}\\
\label{eq:vutfromvin_short}
\end{equation}

\subsection{Utledning av formel for bitsekvenser}
Fra formelene i \cref{eq:vutfromvin_full} og \cref{eq:vutfromvin_short} kan en formel for $v_{ut}$ med $b_0, b_1, b_2$ som ukjente utledes ved hjelp av faktorisering.

\begin{equation}
v_{ut} = \frac{R_LV_H}{2R + R_L}(\frac{b_0}{8} + \frac{b_1}{4} + \frac{b_2}{2})\\
\label{eq:vutfrombin_full}
\end{equation}

\begin{equation}
v_{ut} = V_H(\frac{b_0}{8} + \frac{b_1}{4} + \frac{b_2}{2})\\
\label{eq:vutfrombin_short}
\end{equation}

\Cref{eq:vutfrombin_full} er utledet fra \cref{eq:vutfromvin_full} og \cref{eq:vutfrombin_short} er utledet fra \cref{eq:vutfromvin_short}.

\subsection{Tredje ting}
hensikt

innhold

\section{Praktisk arbeid}

\subsection{Første ting}
\subsection{Andre ting}
\subsection{Tredje ting}

\section{Resultater}

\section{Diskusjon}

blah blah

\section{Konklusjon}

smrt smrt

referansebruk: noe står i læreboka \cite{laerebok}

\begin{thebibliography}{9}
\bibitem{laerebok}
Nilsson og Riedel \textit{Electric Circuits}, 11. Utgave, Pearson.
\bibitem{analogtech}
\url{https://www.analogtechnologies.com/document/white_paper/AWPLED-04-LED-Forward-Voltage.pdf}
\end{thebibliography}

\end{document}
